\newpage
\section*{Lesson 1 Plan}

\textbf{Duration:} 60 minutes

\textbf{Lesson title:} Planning and Prototyping a Meaningful Minecraft Place

\textbf{Student-friendly target:} I can plan and begin building a Minecraft place that shows something important about identity or community, and I can explain how my design choices carry meaning.

\textbf{Identity-safe teacher language:} Your place can be real, imagined, symbolic, collaborative, or future-facing. You do not have to share anything private. The goal is to make thoughtful design choices and explain what those choices communicate.

\textbf{Intellect bridge:} \textcolor{LessonTwoAccent}{Students compare a generic Minecraft structure with a meaningful digital place and ask what makes a digital space readable, welcoming, accessible, or meaningful to an audience. This frames building as intellectual design, not decoration.}

\textbf{Objectives:}
\begin{enumerate}[leftmargin=*]
    \item Students will choose a meaningful place connected to identity, community, belonging, or future possibility.
    \item Students will break the place into smaller build/code steps.
    \item Students will analyze how a sign, path, structure, color, repeated pattern, or coded feature can guide an audience's understanding.
    \item Students will draft the first part of a Developer Design Note by naming their design goal and first 3--5 build/code steps.
\end{enumerate}

\renewcommand{\arraystretch}{0.82}
\setlength{\tabcolsep}{3pt}
\begin{tabularx}{\textwidth}{|>{\raggedright\arraybackslash}p{0.08\textwidth}|Y|Y|}
\hline
\rowcolor{LessonOneSoft}
\textbf{Mins} & \textbf{Activity} & \textbf{Purpose} \\ \hline
0--7 & Compare a generic Minecraft build with a meaningful digital place. Ask what makes a place meaningful, welcoming, or accessible. & Frames design as meaning-making and analysis. \\ \hline
7--15 & Model a small place plan with entrance, sign, path, gathering spot, and repeated pattern. Think aloud about what each feature communicates. & Makes audience and design choices visible. \\ \hline
15--38 & Students choose a real, imagined, symbolic, or future place; sketch 3--5 features; and label first build/code steps. & Builds identity-safe agency and decomposition. \\ \hline
38--50 & Students begin a prototype, test one feature, and add one note explaining what the feature helps the audience understand. & Connects design to audience meaning. \\ \hline
50--60 & Partner share and exit slip: ``This helps my audience understand...'' and ``My next build/code step is...'' & Captures oral and written evidence. \\ \hline
\end{tabularx}

\textbf{Differentiation:} Students may work alone or with roles such as builder, navigator, debugger, or explainer. Beginners may build one strong scene; advanced students may add loops, agent movement, signage, or an audience-centered accessibility feature.

\newpage
\section*{Lesson 2 Plan}

\textbf{Duration:} 60 minutes

\textbf{Lesson title:} Building, Coding, Debugging, and Analyzing the Meaningful Place

\textbf{Student-friendly target:} I can build and revise my Minecraft place so that at least one coded or repeated feature helps communicate meaning.

\textbf{Joy + identity bridge:} Lesson 2 keeps coding joyful by treating debugging as experimentation rather than failure. Students revise features so their worlds become clearer, more welcoming, or more meaningful to an audience.

\textbf{Intellect bridge:} \textcolor{LessonTwoAccent}{Students think like designers and beginning developers. They do not only ask, ``Does it work?'' They also ask, ``What does this design choice help my audience understand?''}

\textbf{Objectives:}
\begin{enumerate}[leftmargin=*]
    \item Students will build at least two meaningful features from their plan.
    \item Students will use a sequence, loop, repeated structure, or agent-based action to create or improve the world.
    \item Students will debug one issue by naming the problem, testing a change, and recording what changed.
    \item Students will analyze how one code/build choice affects audience understanding, welcome, or access.
\end{enumerate}

\renewcommand{\arraystretch}{0.82}
\setlength{\tabcolsep}{3pt}
\begin{tabularx}{\textwidth}{|>{\raggedright\arraybackslash}p{0.08\textwidth}|Y|Y|}
\hline
\rowcolor{LessonTwoSoft}
\textbf{Mins} & \textbf{Activity} & \textbf{Purpose} \\ \hline
0--8 & Students choose one must-build feature and one try-to-code feature, then name what the feature should help the audience understand. & Connects coding/design to meaning. \\ \hline
8--15 & Teacher models a debug cycle and how to write what the revision communicates. & Makes debugging visible, writable, and purposeful. \\ \hline
15--42 & Students build and code with partner roles. Teacher asks: ``What does this choice help your audience read or understand?'' & Develops skill and analysis. \\ \hline
42--52 & Screenshot checkpoint: students label one feature, explain how it works, and explain what it communicates. & Connects product to intellectual explanation. \\ \hline
52--60 & Quick share: one success, one bug, one next step, and one audience meaning. & Builds design language and community expertise. \\ \hline
\end{tabularx}

\textbf{Disciplinary writing evidence:} \textcolor{LessonTwoAccent}{Students complete the middle section of the Developer Design Note: ``My bug or design problem was...'' ``I expected...'' ``What happened was...'' ``I changed...'' ``This helped my audience understand...''}

\newpage
\section*{Lesson 3 Plan}

\textbf{Duration:} 60 minutes

\textbf{Lesson title:} Sharing Worlds, Reading Design, and Reflecting on Representation

\textbf{Student-friendly target:} I can present my Minecraft place, explain how it carries identity or community knowledge, and respond respectfully to other students' designs.

\textbf{Representation bridge:} \textcolor{LessonTwoAccent}{Before students share, the teacher names that digital worlds are designed from choices. Some choices make people feel welcomed and seen; other choices erase people, flatten culture, or make access harder. Students explain what they chose, not private identity.}

\textbf{Objectives:}
\begin{enumerate}[leftmargin=*]
    \item Students will explain at least two design choices and how they carry meaning.
    \item Students will give peer feedback using notice/wonder language.
    \item Students will analyze how digital worlds can include, exclude, or misrepresent people and communities.
    \item Students will revise one feature for clarity, welcome, or accessibility.
\end{enumerate}

\renewcommand{\arraystretch}{0.82}
\setlength{\tabcolsep}{3pt}
\begin{tabularx}{\textwidth}{|>{\raggedright\arraybackslash}p{0.08\textwidth}|Y|p{0.28\textwidth}|Y|}
\hline
\rowcolor{LessonThreeSoft}
\textbf{Mins} & \textbf{Activity} & \textbf{Purpose} \\ \hline
0--8 & Students prepare a two-sentence explanation and choose one screenshot or live view to show. & Supports concise presentation. \\ \hline
8--35 & Gallery walk or paired screen-share. Students explain their worlds and collect one notice/wonder comment. & Centers oral, visual, and digital literacy. \\ \hline
35--48 & Whole-group brainstorm: ``Whose stories do digital worlds usually show, whose could be missing, and how can designers make spaces more respectful and accessible?'' & Keeps Criticality in brainstorm phase while making the representation visible. \\ \hline
48--56 & Students revise one sentence, sign, pathway, or accessibility feature based on feedback and explain why the change matters. & Turns feedback into intellectual action. \\ \hline
56--60 & Exit reflection: ``My world shows...'' ``A designer can make digital spaces more welcoming by...'' ``One feature I revised was...'' & Captures final evidence. \\ \hline
\end{tabularx}

\textbf{Student reflection and analysis prompts:}
\begin{itemize}[leftmargin=*]
    \item \textcolor{LessonTwoAccent}{What does your Minecraft world help an audience understand?}
    \item \textcolor{LessonTwoAccent}{Which design choice carries the most meaning, and why?}
    \item \textcolor{LessonTwoAccent}{How do your paths, signs, colors, structures, repeated patterns, or code guide the audience?}
    \item \textcolor{LessonTwoAccent}{How can a digital world make someone feel welcomed or excluded?}
    \item \textcolor{LessonTwoAccent}{Whose stories are easy to represent in digital worlds, and whose stories are often missing?}
    \item \textcolor{LessonTwoAccent}{How can designers avoid using culture, identity, or community knowledge as decoration?}
    \item \textcolor{LessonTwoAccent}{What did you revise after feedback, and how did that change your world's meaning?}
\end{itemize}

\textbf{Developer Design Note prompt:} \textcolor{LessonTwoAccent}{``My design goal was...'' ``The build/code steps I used were...'' ``One bug or revision I worked through was...'' ``One feature I want my audience to notice is...'' ``This world communicates identity, community, belonging, or future possibility by...'' ``One revision I made for clarity, welcome, or access was...''}

\newpage
\section*{Assessment and Evidence}

\textbf{Evidence collected across the unit:} planning sheet, debug log, screenshots or sketches, teacher conference notes, peer feedback, short presentation, Developer Design Note, and exit reflection.

\renewcommand{\arraystretch}{0.95}
\setlength{\tabcolsep}{3pt}
\begin{tabularx}{\textwidth}{|>{\raggedright\arraybackslash}p{0.17\textwidth}|Y|Y|}
\hline
\rowcolor{SoftBlue}
\textbf{Criterion} & \textbf{Meeting} & \textbf{Developing} \\ \hline
Meaningful design & World includes features connected to identity, community, belonging, or future possibility. & World has features, but the meaning needs clearer explanation. \\ \hline
Computational thinking & Student decomposes, sequences, tests, debugs, or uses a repeated/coded feature. & Student builds but needs support naming steps or debugging. \\ \hline
Disciplinary writing & Student documents design goal, steps, bug/revision, and audience meaning using beginning computer science language. & Student explains the project generally but needs support connecting technical steps to meaning. \\ \hline
Intellectual analysis & \textcolor{LessonTwoAccent}{Student analyzes how design choices such as paths, signs, structures, colors, repeated patterns, code, or accessibility supports help an audience read the world.} & \textcolor{LessonTwoAccent}{Student names a design choice but needs support explaining how it affects audience understanding.} \\ \hline
Critical awareness & Student can brainstorm whose stories are represented, whose stories may be missing, and how design can welcome an audience. & Student notices audience or representation with prompting. \\ \hline
\end{tabularx}

\textcolor{LessonTwoAccent}{\textbf{How the evidence is read:} The teacher does not grade cultural disclosure or artistic polish. The main evidence is whether students connect a design choice to a purpose, use basic computer science language to describe process, analyze how choices affect audience meaning, and revise after testing or feedback.}

\section*{Sample and Appendix Tools}

\textbf{Teacher model sample:} A small community garden world with an entrance path, welcome sign, gathering circle, repeated fence pattern, and one coded/agent-assisted repeated feature. The model should be meaningful but not perfect so students see revision as normal.

\textbf{Student-facing routine:} \textcolor{LessonTwoAccent}{Students use one tool at a time: plan first, debug next, then complete the Developer Design Note and share card. This keeps the workflow accessible while preserving a full evidence trail.}

\newpage
\section*{Reflection}

This mini-unit shows how digital literacy can become a space for belonging. Joy is built into choice, creation, debugging, peer support, and visible progress. Identity is built into meaningful place-making and student control over what is shared. Skills are built through sequencing, decomposition, debugging, documentation, and presentation. \textcolor{LessonTwoAccent}{Intellect is built when students analyze how those technical and design choices communicate meaning to an audience.}

\textcolor{LessonTwoAccent}{The revision strengthens the mini-unit by making intellectual analysis visible. Students are not only building in Minecraft; they are learning how designers and beginning developers communicate a design goal, document a process, test and debug, respond to feedback, and explain how technical choices affect an audience. They learn that code is not only functional. Code, space, signs, paths, colors, structures, repeated patterns, and accessibility supports can all shape how a digital world is read.}

\textcolor{LessonTwoAccent}{For Ethan specifically, the design matters because it does not ask him to simplify himself into one identity label. It gives him a way to connect coding, game design, land, community, mixed identity, and cultural knowledge only to the extent he chooses. That makes the technical work more humane: code becomes a tool for meaning-making, not a gatekeeping measure.}

For my own teaching, the most important design decision is that students are assessed on how intentionally they communicate meaning. That distinction matters for neurodivergent, culturally connected, multilingual, and mixed-identity students because it keeps literacy tied to agency. In this unit, joy is not separate from rigor: students still decompose, debug, revise, document, analyze, and explain, but those skills serve a world they care about.

\textcolor{LessonTwoAccent}{This pursuit expands students' intellectualism because it helps them understand that design is a form of thinking. Students learn to read digital spaces as texts and to understand that technical choices are not neutral. A design can welcome, confuse, include, exclude, preserve, flatten, or erase meaning. The work prepares students to see themselves not just as users of technology, but as designers whose choices affect others.}

\newpage
